\section{Introduction and objective}

Dense matrix multiplication is the standard microbenchmark for studying how a programming language
implementation turns the same algorithm into machine work. This study compares the \emph{textbook triple-loop}
algorithm, which performs $n^3$ multiply--add operations on $n\times n$ matrices, in three languages:
C (ahead-of-time compiled), Java (JIT-compiled on the JVM) and Python (CPython, interpreted).

\paragraph{Research question.} \emph{How do the runtime and the memory use of the basic $O(n^3)$ multiplication of
$n\times n$ \texttt{float64} matrices grow with $n$ in C, Java and Python, when the three programs use the same
inputs, the same loop order and the same arithmetic on one documented machine?}

\paragraph{Factor and outcomes.} The factor that changes is the matrix size $n$ (from $10$ up to the largest size
that fits the declared budgets); the outcomes are the time of one multiplication (kernel only) and the peak resident
memory of the process.

\paragraph{Hypotheses (predictions, stated before running the experiments).}
\begin{enumerate}
  \item[H1] In every language the time grows approximately as $n^3$ (log--log slope $\approx 3$), but the constant
        factor differs.
  \item[H2] The order of speed is C $\lesssim$ Java $\ll$ Python, with Python one to two orders of magnitude slower
        than the compiled languages, because of interpreter dispatch and boxed floating-point objects.
  \item[H3] The cost per multiply--add is not constant: once the matrix $B$, which is read by columns, no longer
        fits in a cache level, the cost per operation increases.
  \item[H4] Memory is dominated by array storage: about $24\,n^2$ bytes for C and Java ($3$ matrices $\times 8$ bytes), a
        much larger $\approx 120\,n^2$ bytes for Python (each element is an 8-byte pointer plus a separately allocated
        float object), plus a fixed runtime baseline that is largest for the JVM.
\end{enumerate}

\paragraph{Scope.} Conclusions are limited to the implementations, compiler and runtime settings, inputs and the
virtual machine described in Section~\ref{sec:env}. They are \emph{not} a general ranking of programming languages.
Language-specific optimisations (loop reordering, blocking, libraries) are deliberately excluded; they belong to
the later assignments.

\section{Background}

\paragraph{Algorithm.} For $C = AB$ with $A,B\in\mathbb{R}^{n\times n}$, $c_{ij}=\sum_{k=1}^{n}a_{ik}b_{kj}$. The
direct evaluation needs $n^3$ multiplications and $n^3-n^2$ additions, hence $\Theta(n^3)$ work \cite{clrs}. All
three programs use the loop order $i$--$j$--$k$ and accumulate each $c_{ij}$ in a scalar before storing it.

\paragraph{Memory access pattern.} With row-major storage the inner loop walks $a_{ik}$ along a row (unit stride) but
$b_{kj}$ down a column (stride $8n$ bytes). Once a column of $B$ touches more cache lines than the cache can
keep for the next value of $j$, every access to $B$ becomes a cache miss; for stride $8n$ with $n$ a power of two,
many addresses additionally map to the same cache sets \cite{drepper,hennessy}. This is why the cost per
multiply--add is expected to rise with $n$ even though the operation count is exactly cubic.

\paragraph{Execution models.} C code is compiled to machine code before running. Java bytecode starts interpreted
and is compiled by HotSpot's tiered JIT (C1, then C2), including on-stack replacement of hot loops; array accesses
keep bounds checks unless the JIT can remove them. CPython executes bytecode in an interpreter: every
\texttt{A[i][k]*B[k][j]} involves several bytecode dispatches, index operations and the allocation of a new
\texttt{float} object.

\paragraph{Floating-point error.} For a sum of $n$ products, any evaluation order gives
$|\hat c-c|\le\gamma_n\sum_k|a_{ik}b_{kj}|$ with $\gamma_n=nu/(1-nu)$ and $u=2^{-53}$ \cite{higham}. Because all
entries here are non-negative, $\sum_k|a_{ik}b_{kj}|=c$, so each implementation has relative error at most
$\gamma_n$, and two different evaluation orders differ by at most $2\gamma_n$ \cite{goldberg}. This bound is used as the
acceptance tolerance (Section~\ref{sec:correct}).

\paragraph{Benchmarking practice.} The protocol follows common recommendations: report central value and spread
over repeated runs, separate warm-up from steady state for JIT-compiled languages, state exactly what is timed
and report the environment \cite{georges,hoefler}.

\section{Methodology}

\subsection{Implementations}
\label{sec:impl}
\begin{itemize}
  \item \textbf{C} (\texttt{src/c/matmul.c}): one contiguous row-major block of $n^2$ \texttt{double} per matrix; the
        kernel is a separate non-inlined function so repetitions cannot be merged by the compiler.
  \item \textbf{Java} (\texttt{src/java/MatMul.java}): \texttt{double[n][n]}, i.e.\ $n$ separately allocated row
        arrays; kernel in its own static method.
  \item \textbf{Python} (\texttt{src/python/matmul.py}): list of $n$ lists of $n$ boxed \texttt{float}; pure Python
        loops, \emph{no} NumPy in the kernel.
\end{itemize}
The kernel is identical in structure in all three:
\texttt{s = 0.0; for k: s += A[i][k]*B[k][j]; C[i][j] = s}, for $i$, $j$ in this order. The scalar accumulator
differs slightly from a \texttt{C[i][j] += \dots} form: it keeps the same operation order (so the results are
bit-for-bit identical, see Section~\ref{sec:correct}) and avoids language-specific aliasing effects in the
comparison. Unavoidable differences are the storage layouts above and the runtime behaviour (JIT, interpreter,
boxing). No library multiplication is used by any implementation.

\subsection{Inputs}
Matrices $A$ and $B$ for each $n$ are generated once with NumPy's \texttt{default\_rng} (PCG64) seeded with
$[20261012, n, w]$ ($w=0$ for $A$, $w=1$ for $B$), uniform in $[0,1)$, and stored as raw little-endian binary64
files. All three programs read the \emph{same files}, so the logical inputs are identical and no cross-language
random generator is involved. SHA-256 digests of all files are stored in \texttt{data/manifest.json}.

\subsection{Environment}
\label{sec:env}
All measurements were taken on one machine, sequentially, with no other workload running (the driver only polled
the child process every $0.5$\,s). Table~\ref{tab:env} lists the environment. The machine is a virtual machine
(KVM) with $2$ vCPUs; the frequency governor is not visible and the processes were not pinned to a CPU. Each
kernel is single-threaded; the JVM may place its JIT and GC helper threads on the second vCPU. Thread counts used
by numerical libraries are not relevant because no library is used in the kernels (NumPy is only used for
generating inputs and checking results, \emph{outside} the measurements) \cite{numpy}.

\begin{table}[H]
\centering\small
\begin{tabular}{lp{10.5cm}}
\toprule
Item & Value \\
\midrule
CPU & Intel(R) Xeon(R) Processor @ 2.80GHz \\
Cores / logical CPUs & 2 physical (vCPU), 2 logical; 1 thread(s)/core \\
Caches & L1d 64 KiB (2 instances), L2 2 MiB (2 instances), L3 33 MiB (1 instance) \\
RAM & 8030 MiB, swap 0 MiB \\
Virtualisation & KVM (full) \\
OS / kernel & Ubuntu 24.04.5 LTS, Linux 6.18.44-fc-v114 \\
Python & CPython 3.13.16 (GCC 13.3.0) \\
Java & OpenJDK 21.0.12.1 (HotSpot 64-bit server VM); flags -XX:+UseSerialGC -Xms256m -Xmx4g (default tiered JIT: C1+C2) \\
C compiler & gcc (Ubuntu 13.3.0-6ubuntu2~24.04.1) 13.3.0; flags -O2 -std=c11 -ffp-contract=off  (no -march=native; no -ffast-math) \\
NumPy (reference only) & 2.5.3 \\
Threads in the kernel & 1 (single-threaded); no CPU pinning; frequency governor: not exposed (virtual machine) \\
\bottomrule
\end{tabular}

\caption{Hardware and software environment.}
\label{tab:env}
\end{table}

\subsection{Fixed settings}
Fixed in every comparison: algorithm and loop order, numeric type (binary64), inputs, single thread, one process
per configuration, monotonic timers, kernel-only timing. C is compiled with
\texttt{gcc -O2 -std=c11 -ffp-contract=off} (no \texttt{-march=native}, no \texttt{-ffast-math}): this forbids fused
multiply--add and re-association, so the arithmetic is the same as in Java and Python. Java runs on OpenJDK~21
with the default tiered JIT and \texttt{-XX:+UseSerialGC -Xms256m -Xmx4g}. Python is CPython~3.13 with the cyclic
garbage collector disabled inside the timed region (the kernel allocates only floats, which are not GC-tracked, so
this only removes a source of pauses). Different compiler options for C are explored in a \emph{separate}
supplementary experiment (Section~\ref{sec:cflags}).

\subsection{What the timer includes}
\label{sec:timer}
\textbf{Kernel only.} The timed region is the triple loop. Reading the inputs, allocating $C$, zero-initialising
$C$ (so its pages are touched before timing), computing the checksum and dumping the result are \emph{outside} the
timer. In Python the row lists of $C$ are created outside the timer; the float objects that hold the results are
created by the kernel and are therefore inside it. Timers: \texttt{clock\_gettime(CLOCK\_MONOTONIC)} (C),
\texttt{System.nanoTime()} (Java), \texttt{time.perf\_counter\_ns()} (Python). The result of every multiplication is
consumed: after each timed sample a checksum of $C$ (plain left-to-right sum in row-major order) is computed and
recorded, and the final $C$ is written to disk and verified.

\subsection{Repetition, warm-up and summary statistics}
Each configuration $(\text{language},n)$ runs in its own process with a number of \emph{warm-up} samples
(excluded from the statistics but kept in the raw data) followed by \emph{measured} samples. The numbers depend
on the estimated time of one multiplication $\hat t$: $\hat t\le 1$\,s: 3 warm-up, 10 measured; $\le30$\,s: 1 warm-up,
5 measured; $\le120$\,s: no warm-up, 3 measured; larger: no warm-up, 2 measured. Fewer repetitions for costly runs
are a compromise forced by the budget (a single Python multiplication at the largest sizes takes minutes); the
resulting spread is reported through min--max and IQR, and the number of repetitions is given for every
configuration (Table~\ref{tab:time}). When $\hat t<50$\,ms a batch of $m=\lceil 0.1\,\text{s}/\hat t\,\rceil$
multiplications is timed and divided by $m$, and the reported value is the \emph{time per multiplication}. Reported
statistics are the \textbf{median} and the \textbf{inter-quartile range} (IQR, $Q_3-Q_1$). The runs for different
languages are interleaved by size (for each $n$: C, Java, Python) so that slow drift of the machine affects all
languages alike.

\subsection{Budgets and limits}
\label{sec:budgets}
Declared before the large runs: a single timed sample must complete within \textbf{300\,s}; peak resident memory per
run must not exceed \textbf{4\,GiB}. A run that violates a limit is killed and recorded as a \emph{timeout} or
\emph{memory failure}, never given an invented time. The driver does not attempt sizes whose predicted time exceeds
three times the limit and stops enlarging $n$ for an implementation after its first failure. The \emph{largest
feasible} size (all repetitions completed, results validated) and the \emph{largest attempted} size are reported
separately (Table~\ref{tab:limits}). Planning aid: three binary64 $n\times n$ matrices need $24n^2$ bytes
(e.g.\ $\approx 0.4$\,GB at $n=4096$), so memory was not expected to bind before time did.

\subsection{Memory metric}
The primary metric is the \textbf{peak resident set size} of the whole process (\texttt{VmHWM} from
\texttt{/proc/self/status}, read by each program just before it exits and before the result is dumped), so it
includes the runtime itself (JVM, interpreter) and all matrices. The kernel's own accounting
(\texttt{ru\_maxrss} from \texttt{wait4}) is stored as a cross-check. RSS is not the same quantity as managed-heap use
or as the arithmetic estimate of array storage, and it is not comparable with them directly; the model
lines in Figure~\ref{fig:mem} are estimates of array storage added to the measured baseline of the smallest run.
\textbf{Limitations:} RSS counts pages actually resident (not reserved), includes loader/library pages, and has
granularity of one page; the Python process also needs transient buffers while loading rows.

\subsection{Correctness}
\label{sec:correct}
Before any benchmark, \texttt{scripts/test\_correctness.py} runs all three programs on small cases whose answers are
known (scalar, hand-computed $3\times3$, non-commutative $2\times2$, zero matrix, identity on both sides,
all-ones, small-integer matrices; \emph{exact} equality required, all intermediate values are integers below $2^{53}$,
so there is no overflow or rounding), and on random matrices against NumPy's \texttt{float64} product (tolerance
$2\gamma_n$ relative, derived in Section~2; absolute tolerance $0$ because every reference entry is positive). It also
requires the three languages to produce \emph{bit-for-bit identical} results. During the benchmark every run dumps
its final product and the driver checks it element-wise against NumPy with the same tolerance, plus the checksum
against the reference sum (tolerance $\gamma_{n^2+2n}$). NumPy/BLAS serves only as a trusted reference and is
not one of the implementations under study.

\subsection{Supplementary experiments (exploratory, kept separate)}
\label{sec:supp}
The following experiments were designed \emph{after} looking at the main results, to test explanations for
patterns seen there; they are therefore exploratory, they are reported separately, and they do not change the
controlled comparison above. All use the same driver, inputs, budgets and validation.
\begin{enumerate}
  \item \textbf{C optimisation level} (Section~\ref{sec:cflags}): the C program compiled with \texttt{-O0}, \texttt{-O1},
        \texttt{-O2}, \texttt{-O3} and \texttt{-O3 -march=native} (always with \texttt{-ffp-contract=off}) at
        $n\in\{512,1000,1024\}$.
  \item \textbf{Sizes around powers of two} (Section~\ref{sec:pow2}): $n\in\{496,504,512,520\}$ and
        $\{1016,1024,1032\}$ for C, Java, and the variant below.
  \item \textbf{Java with a flat array} (\texttt{src/java/MatMulFlat.java}): the same Java algorithm with each matrix
        stored as one row-major \texttt{double[n*n]} (the layout used by the C program), to test whether the
        jagged \texttt{double[n][n]} layout explains the Java--C gap. Sizes $n\in\{500,1000,1250,1500,2000\}$ and the
        power-of-two neighbourhoods. Correctness is checked by the driver on every run (same tolerance), not by
        the unit-test script.
  \item \textbf{Between-process variability} (Section~\ref{sec:procvar}): C and Java at $n=1024$ in six
        independent processes (the main run, the power-of-two run and four dedicated repeats).
\end{enumerate}
